\begin{frame}{Funções de transferência com componentes medidos}
  \small
  \begin{columns}[T,onlytextwidth]
    \begin{column}{0.31\textwidth}
      \begin{block}{1. Sem carga}
        \[
          H_1(s)\approx\frac{1}{1+0{,}00226710s}.
        \]
        \[
          \tau_1=2{,}2671\,\mathrm{ms}
        \]
        \[
          K_1=1
        \]
        \[
          f_{c1}\approx70{,}20\,\mathrm{Hz}.
        \]
      \end{block}
    \end{column}
    \begin{column}{0.31\textwidth}
      \begin{block}{2. Carga direta}
        \[
          H_2(s)\approx\frac{0{,}59472}{1+0{,}00145238s}.
        \]
        \[
          \tau_2=1{,}4524\,\mathrm{ms}
        \]
        \[
          K_2\approx0{,}59472
        \]
        \[
          f_{c2}\approx109{,}58\,\mathrm{Hz}.
        \]
      \end{block}
    \end{column}
    \begin{column}{0.31\textwidth}
      \begin{block}{3. Com buffer}
        \[
          H_3(s)\approx\frac{1}{1+0{,}00231233s}.
        \]
        \[
          \tau_3=2{,}3123\,\mathrm{ms}
        \]
        \[
          K_3=1
        \]
        \[
          f_{c3}\approx68{,}83\,\mathrm{Hz}.
        \]
      \end{block}
    \end{column}
  \end{columns}

  \vspace{0.12cm}
  \centering\scriptsize
  No circuito 2, $R\parallel R_L\approx13{,}0963\,\mathrm{k}\Omega$.
  As diferenças entre os cortes dos circuitos 1 e 3 decorrem dos componentes medidos.
\end{frame}
